\section{Introduction}
\label{sec:intro}

Modern multimodal models are increasingly evaluated on their ability to turn
images into code. Work on screenshot-to-HTML~\cite{beltramelli2017pix2code,si2024design2code},
structure extraction from webpages, LaTeX, and music
scores~\cite{roberts2024image2struct}, and symbolic vector
generation~\cite{lin2025vcode,zhou2026omni} all ask the same underlying question
in different clothes: \emph{can a model look at a picture and produce the program
that generated it?} Earlier work on visual program
induction~\cite{sharma2018csgnet,liu2019scenes,li2020perspective,li2020multiplane,duan2022parametric,grand2024lilo}
framed this problem as inverse graphics over a small symbolic DSL, and more
recently \textsc{TurtleBench}~\cite{rismanchian2025turtlebench} has revived it
as a benchmark-first target for large vision-language models.

Across these lines of work, three design pressures keep recurring: (1) deterministic
execution so that scoring is principled; (2) render-based scoring so that
semantically equivalent programs are not penalized for textual
differences~\cite{roberts2024image2struct}; and (3) controlled generation so that
failure modes can be attributed to specific visual factors, in the tradition of
\textsc{CLEVR}~\cite{johnson2017clevr,bahdanau2019closure}. Most existing
benchmarks satisfy one or two of these, but few satisfy all three while remaining
\emph{renewable} -- that is, cheap enough to regenerate when an existing split
becomes contaminated. For model development, renewability also changes the
feedback cycle: a researcher can generate fresh instances, run a model, and
obtain objective scores without commissioning new labels or manual judgments for
each refreshed example.

We present \textsc{ShapeCodeBench}, a perception-to-program benchmark that attempts to
hit the intersection of these pressures. The task is narrow by design: the DSL
has exactly four primitives (\texttt{filled\_circle}, \texttt{circle},
\texttt{filled\_square}, \texttt{square}) and the canvas is fixed at
$512{\times}512$ grayscale with a black foreground on a white background. Every
sample is generated from a seeded RNG with explicit difficulty controls on shape
count, size, stroke width, overlap, and canvas clipping. The evaluator parses
predictions with a safe restricted Python parser, re-renders them through the
same deterministic Pillow pipeline used to produce the target, and compares
rasters.

\paragraph{Contributions.} We make the following contributions:
\begin{enumerate}
  \item We release the \textsc{ShapeCodeBench} benchmark: a four-primitive drawing DSL,
    a safe restricted parser, a seeded scene generator with three difficulty
    tiers, and a render-based evaluator with five primary metrics.
  \item We freeze an evaluation split, \textsc{eval\_v1}, of 150 samples with
    deterministic seeds, and publish per-sample raster hashes to make the exact
    evaluation instances reproducible across platforms.
  \item We make benchmark refresh a first-class workflow: new held-out splits can
    be generated from fresh seeds and scored automatically, enabling fast
    regression-style feedback while avoiding per-instance human annotation or
    manual judging. This mitigates exact-instance contamination but does not
    prevent models from learning the generator distribution.
  \item We release a provider-agnostic runner that records prompts, model
    configuration, raw outputs, normalized predictions, metrics, and per-sample
    artifacts, making model evaluations auditable and easy to extend.
  \item We report baseline results for four multimodal model configurations
    -- \textsc{Claude Opus 4.7} (1M context) at \texttt{high} and \texttt{max}
    effort, and \textsc{GPT-5.5} at \texttt{medium} and \texttt{extra\_high}
    reasoning effort -- alongside two non-LLM baselines (empty program,
    classical-CV heuristic). The strongest multimodal model by foreground IoU
    (\textsc{GPT-5.5}/\texttt{extra\_high}) reaches mean foreground IoU $0.87$,
    the best multimodal exact match remains $0.027$, and the classical
    heuristic leads easy-tier exact match at $0.26$ -- a
    tier-dependent crossover that confirms \textsc{ShapeCodeBench} is not saturated
    and exposes distinct failure modes in perception versus structured code
    emission. Effort tier helps modestly (\texttt{max}\,$>$\,\texttt{high} for
    Claude on FG-IoU; \texttt{extra\_high}\,$>$\,\texttt{medium} for GPT-5.5)
    but does not close either gap.
\end{enumerate}

The rest of the paper is organized as follows. Section~\ref{sec:related} places
\textsc{ShapeCodeBench} against prior work. Section~\ref{sec:benchmark} describes the
DSL, generator, renderer, and evaluator. Section~\ref{sec:experiments} details
the experimental setup and headline results. Section~\ref{sec:analysis} analyzes
failure modes, the heuristic-vs-LLM gap, and difficulty validity.
Section~\ref{sec:limitations} discusses limitations and future work.
